\section{PENDAHULUAN}

\subsection{Analisis Situasi}

Perkembangan teknologi kecerdasan buatan atau \textit{Artificial Intelligence} dan \textit{Machine Learning} telah mendorong perubahan signifikan dalam kebutuhan kompetensi sumber daya manusia di bidang teknologi informasi. Industri digital saat ini tidak hanya membutuhkan lulusan yang mampu membuat program komputer, tetapi juga tenaga kerja yang memiliki kemampuan mengolah data, membangun model prediksi, memahami pola dari citra atau teks, serta menerapkan solusi berbasis kecerdasan buatan pada permasalahan nyata. Kondisi ini menjadi tantangan bagi pendidikan vokasional, khususnya Sekolah Menengah Kejuruan, untuk menyesuaikan kompetensi siswa dengan kebutuhan industri.

SMKN 2 Cimahi merupakan salah satu sekolah kejuruan yang memiliki Jurusan Rekayasa Perangkat Lunak dan menjadi mitra dalam kegiatan Pengabdian kepada Masyarakat ini. Jurusan Rekayasa Perangkat Lunak memiliki karakteristik pembelajaran yang dekat dengan pemrograman, pengembangan aplikasi, basis data, dan praktik rekayasa perangkat lunak. Berdasarkan rancangan kegiatan, sasaran program ini adalah siswa Jurusan Rekayasa Perangkat Lunak SMKN 2 Cimahi yang telah memiliki dasar pemrograman. Selain itu, sekolah mitra juga telah memiliki hubungan kemitraan sebelumnya dengan ULBI melalui kegiatan PKM Workshop Bootcamp Laravel pada tahun 2025, sehingga kegiatan ini menjadi bentuk keberlanjutan penguatan kompetensi teknologi bagi siswa SMK.

Pada kegiatan PKM sebelumnya, penguatan kompetensi siswa diarahkan pada pengembangan aplikasi web menggunakan Laravel melalui pendekatan \textit{Project-Based Learning}. Kegiatan tersebut menunjukkan bahwa siswa SMK dapat diarahkan untuk menghasilkan produk teknologi sederhana apabila diberikan modul, pendampingan, dan studi kasus yang sesuai dengan tingkat kemampuan mereka. Pola ini menjadi dasar penting dalam merancang kegiatan lanjutan yang lebih relevan dengan perkembangan teknologi terkini, yaitu pengenalan \textit{Machine Learning} dan \textit{Computer Vision} melalui studi kasus yang sederhana, visual, dan aplikatif.

Berdasarkan identifikasi awal, siswa Jurusan Rekayasa Perangkat Lunak SMKN 2 Cimahi telah memiliki fondasi pemrograman dasar, logika pemrograman, serta pengalaman awal dalam pengembangan aplikasi. Namun, pembelajaran yang diterima siswa belum secara khusus mengarah pada praktik \textit{Machine Learning} dan \textit{Computer Vision} secara terstruktur. Siswa belum memiliki pengalaman membangun alur kerja \textit{Machine Learning} secara \textit{end-to-end}, mulai dari pengenalan dataset, praproses data, pelatihan model, evaluasi hasil, hingga implementasi model ke dalam aplikasi sederhana. Kondisi ini menunjukkan adanya kesenjangan antara kemampuan pemrograman dasar yang sudah dimiliki siswa dengan kebutuhan kompetensi kecerdasan buatan yang semakin dibutuhkan dalam dunia kerja teknologi.

Kesenjangan tersebut menjadi peluang bagi ULBI untuk melakukan transfer ilmu pengetahuan dan teknologi melalui kegiatan \textit{Workshop Machine Learning} untuk Prediksi Awal Kesehatan Tanaman melalui Analisis Citra Daun. Studi kasus citra daun dipilih karena mudah divisualisasikan, dapat dipahami oleh siswa, menggunakan dataset publik yang tersedia secara gratis, serta tidak memerlukan perangkat keras khusus. Melalui kegiatan ini, siswa tidak hanya mempelajari konsep dasar \textit{Machine Learning}, tetapi juga memperoleh pengalaman praktik dalam membangun model klasifikasi citra dan aplikasi sederhana yang dapat digunakan sebagai portofolio pembelajaran.

Pelatihan \textit{Machine Learning} berbasis proyek ini diharapkan dapat memberikan manfaat langsung bagi mitra. Bagi siswa, kegiatan ini dapat meningkatkan pemahaman terhadap penerapan kecerdasan buatan, memperkaya portofolio proyek teknologi, dan meningkatkan kesiapan menghadapi dunia kerja maupun studi lanjut di bidang teknologi informasi. Bagi sekolah, kegiatan ini dapat menghasilkan modul pelatihan dan bahan ajar praktis yang dapat digunakan kembali dalam pembelajaran. Dengan demikian, kegiatan PKM ini tidak hanya berorientasi pada pelatihan sesaat, tetapi juga diarahkan untuk mendukung penguatan kompetensi digital siswa SMK secara berkelanjutan.

\subsection{Permasalahan Mitra}

Berdasarkan analisis situasi yang telah diuraikan pada sub bab 1.1, beberapa permasalahan prioritas yang dihadapi mitra dan perlu diselesaikan melalui kegiatan Pengabdian kepada Masyarakat ini adalah sebagai berikut:

\begin{enumerate}[leftmargin=*, label=\arabic*.]
    \item \textbf{Keterbatasan materi praktis \textit{Machine Learning} dan \textit{Computer Vision}}

    Materi pembelajaran yang diterima siswa belum secara khusus mengarah pada praktik \textit{Machine Learning} dan \textit{Computer Vision}, sehingga siswa belum memperoleh pengalaman nyata dalam membangun model kecerdasan buatan yang sesuai dengan perkembangan industri digital.

    \item \textbf{Minimnya akses terhadap \textit{tools}, dataset, dan referensi pembelajaran}

    Pembelajaran \textit{Machine Learning} sering dianggap sulit karena membutuhkan pemahaman matematika, statistik, dan konfigurasi perangkat lunak tertentu. Siswa membutuhkan media pembelajaran yang lebih sederhana dan dapat diakses tanpa instalasi rumit, sehingga penggunaan \textit{tools} berbasis \textit{browser} seperti \textit{Google Colab} dan \textit{dataset} publik menjadi kebutuhan penting.

    \item \textbf{Belum adanya pengalaman membangun \textit{pipeline Machine Learning} secara \textit{end-to-end}}

    Siswa belum memiliki pengalaman membangun alur kerja \textit{Machine Learning} secara \textit{end-to-end}, mulai dari pengenalan data, praproses citra, pelatihan model, evaluasi, hingga implementasi ke dalam aplikasi sederhana. Pendekatan \textit{Project-Based Learning} dinilai relevan karena mendorong siswa memahami teori sekaligus menerapkannya melalui proyek yang membangun keterampilan praktis dan pemecahan masalah.

    \item \textbf{Belum tersedianya portofolio siswa berbasis kecerdasan buatan}

    Sebagian besar siswa belum memiliki portofolio berbasis \textit{Artificial Intelligence} atau \textit{Machine Learning} yang dapat mendukung kesiapan kerja, magang industri, maupun pendaftaran ke perguruan tinggi.

    \item \textbf{Keterbatasan pemahaman terhadap penerapan \textit{Machine Learning} pada masalah nyata}

    Siswa membutuhkan contoh penerapan \textit{Machine Learning} yang dekat dengan kehidupan sehari-hari agar konsep kecerdasan buatan tidak dipahami sebagai materi abstrak. Studi kasus prediksi awal kesehatan tanaman melalui analisis citra daun dipilih karena bersifat visual dan mudah dipahami, sekaligus melatih siswa memahami manfaat dan keterbatasan model sebagai alat \textit{screening} awal, bukan diagnosis final.
\end{enumerate}
